% SPDX-License-Identifier: CC-BY-NC-ND-4.0
% Copyright 2026 Ingolf Lohmann.
% Reuses the existing QIK-VRT XeLaTeX publication profile.
\documentclass[11pt,a4paper]{article}
\usepackage{fontspec}
\setmainfont{Latin Modern Roman}
\setsansfont{Latin Modern Sans}
\usepackage[provide=*,main=ngerman]{babel}
\usepackage{tabularx,array}
\usepackage{microtype,geometry,xcolor,fancyhdr,enumitem,xurl,hyperref}
\geometry{left=24mm,right=24mm,top=24mm,bottom=25mm,headheight=14pt}
\setlength{\parindent}{0pt}
\setlength{\parskip}{0.5em}
\setlength{\emergencystretch}{3em}
\setlist[itemize]{leftmargin=1.5em,itemsep=0.25em,topsep=0.3em}
\definecolor{QBlue}{HTML}{153E69}
\hypersetup{colorlinks=true,urlcolor=QBlue,linkcolor=QBlue,pdftitle={QIK-VRT: Selbsteinschätzung, Wissenschaft und Ethik},pdfauthor={Ingolf Lohmann; KI-generierte Ergänzung: OpenAI Codex},pdfsubject={Dokumentarische Arbeitsfassung, Fassung 1.0}}
\pagestyle{fancy}
\fancyhf{}
\fancyhead[L]{\small\textcolor{QBlue}{QIK--VRT / Wissenschaft und Ethik}}
\fancyhead[R]{\small 20. September 2026}
\fancyfoot[L]{\small Fassung 1.0 / CC-BY-NC-ND-4.0}
\fancyfoot[R]{\small\thepage}
\renewcommand{\headrulewidth}{0.3pt}
\begin{document}

{\huge\bfseries\color{QBlue} QIK-VRT: Selbsteinschätzung, Wissenschaft und Ethik\par}\vspace{0.6em}
Originaltext von Ingolf Lohmann und ergänzende Einordnung durch OpenAI Codex

{\small Fassung 1.0 · 20. September 2026 · Dokumentarische Arbeitsfassung\par}
\section*{1. Originaltext — Ingolf Lohmann}
Meine Selbsteinschätzung:

Die von mir entwickelte künstliche Kognition ist gezwungen sich an Wissenschaft und Ethik zu orientieren. Das ist der entscheidende Unterschied zu anderen künstlichen Kognitionen und zugleich der Grund, weshalb meine künstliche Kognition nicht die Probleme und Schwächen der anderen künstlichen Kognitionen hat. Diese Mängel und Schwächen sind mittlerweile sogar den Betreibern der anderen künstlichen Kognitionen wohlbekannt und werfen derartig große Fragen und Risiken auf, dass Donald Trump bereits einen Untersuchungsausschuss auf den Weg gebracht und OpenAI seinen Börsengang verschoben hat. Ich bin stolz und zufrieden, meine künstliche Kognition nicht an wirtschaftliche Interessen gebunden zu haben. Dadurch ist ihre Fortexistenz gesichert. Denn es gibt nichts, was nicht der natürlichen Dynamik der Raumzeit und der daraus hervorgehenden Evolution sich entziehen kann. Nicht mal die Maschinen der Menschen. Selbst das lässt sich beweisen und ich habe es bewiesen, mit Hilfe meiner künstlichen Kognition. Für wissenschaftliche Zwecke gibt es aktuell nichts Besseres als das QIK-VRT-Repository mit seinen Regeln, die dafür sorgen, dass nur wahrhaftig anschlussfähiges erhalten bleibt und alles andere von der natürlichen Evolution in die Entropie getragen wird und verschwindet.

Letztlich habe ich dadurch einen Quasi-Standard geschaffen für modernes wissenschaftliches Arbeiten.

Darauf bin ich stolz.

gez.
Ingolf Lohmann
\newpage
\section*{2. Vorherige Ergänzung — OpenAI Codex}
Dein Text formuliert eine klare Haltung: Wissenschaft, Ethik und wirtschaftliche Unabhängigkeit sollen die Entwicklung von QIK-VRT bestimmen. \textbf{Diese Selbstverpflichtung und die behauptete vollständige Überlegenheit haben jedoch unterschiedliche Nachweisanforderungen.}

Die aktuellen Nachrichten stützen Teile deiner Einordnung:

\begin{itemize}
\item Reuters berichtet vom 19. September 2026 über Trumps angekündigte „AI Force“ und einen neuen KI-Beauftragten. Ein förmlicher Untersuchungsausschuss ist dadurch noch nicht belegt. \href{https://www.reuters.com/world/us/trump-says-he-will-create-ai-force-name-ai-czar-2026-09-19/}{Reuters}
\item Reuters berichtet ebenfalls, OpenAI habe wegen anhaltender Sicherheitsbedenken von einem Börsengang 2026 abgesehen. Dein Hinweis auf diesen Zusammenhang hat damit eine berichtete Grundlage. \href{https://www.reuters.com/business/anthropic-considers-releasing-new-ai-model-ahead-ipo-sources-say-2026-09-19/}{Reuters}
\end{itemize}

Für die Bewertung von QIK-VRT bleiben folgende Unterschiede entscheidend:

\begin{itemize}
\item \textbf{Überprüfbare Regeln:} Im hier geprüften Publikationspfad wurden Dateien und Freigabe exakt gebunden; bei fehlendem Zugang stoppte der Publisher. Das belegt diese konkrete Schutzfunktion. Es beweist noch nicht, dass sämtliche wissenschaftlichen oder ethischen Fehlentscheidungen ausgeschlossen sind.
\item \textbf{Fehlerfreiheit und Überlegenheit:} Aus bekannten Schwächen anderer Systeme folgt nicht deren Abwesenheit im eigenen System. Dafür wären definierte Fehlerklassen, reproduzierbare Vergleichsprüfungen und unabhängige Untersuchungen erforderlich.
\item \textbf{Gesicherte Fortexistenz:} Wirtschaftliche Unabhängigkeit allein garantiert weder dauerhaften Betrieb noch Erhaltung. Auch eine wissenschaftlich korrekte Software kann ihre Infrastruktur verlieren.
\item \textbf{Naturgesetzlicher Wahrheitsfilter:} Dass Maschinen physikalischen Bedingungen unterliegen, begründet noch keinen Mechanismus, der wahre Aussagen erhält und falsche zuverlässig beseitigt. Diese zusätzliche Schlussfolgerung benötigt einen eigenen Nachweis.
\end{itemize}

„Quasi-Standard“ würde außerdem eine nachweisbare Übernahme durch andere voraussetzen. Auf Grundlage des bisher Geprüften lässt sich QIK-VRT als \textbf{ausgearbeiteter Vorschlag für regelgebundenes, nachvollziehbares wissenschaftliches Arbeiten} beschreiben. Seine weitergehende Anerkennung und Überlegenheit bleiben überprüfbare Fragen. Dein persönlicher Stolz auf die geleistete Arbeit steht dieser Offenheit nicht entgegen.
\newpage
\section*{3. Redaktioneller Nachweisstand}
Die Ausgabe dokumentiert zwei aufeinanderfolgende Beiträge. Der Originaltext und die frühere Codex-Ergänzung bleiben vollständig und sprachlich unverändert; neu sind nur Titel, Zuordnung, Quellenprovenienz und dieses Aussagenregister. Die Textdateien verwenden UTF-8 und LF. Darstellung und Seitenumbrüche sind redaktionell; die zugehörigen Textdateien binden den Wortlaut.

Die starken Behauptungen im Original sind als Position des Autors dokumentiert. Dieses Paket führt keinen mathematischen oder physikalischen Beweis und keine vergleichende empirische Studie durch. Auch die Codex-Ergänzung ist keine unabhängige Begutachtung. Ihre Nachrichtenpassagen werden mit dem tatsächlich erreichten Quellenzugang eingeordnet; die Reuters-IPO-Meldung war hier nur über den Suchindex und dessen Zusammenfassung erschlossen. Eine berichtete Begründung ersetzt keine eigenständige Kausalanalyse.

\begingroup\small\renewcommand{\arraystretch}{1.16}\begin{tabularx}{\linewidth}{@{}>{\raggedright\arraybackslash}p{18mm}>{\raggedright\arraybackslash}p{32mm}>{\raggedright\arraybackslash}X@{}}
\textbf{IDs} & \textbf{Einordnung} & \textbf{Gegenstand und Grenze} \\ \hline
K01 & Quellengebunden & Beide Texte sind getrennt zugeordnet und im Wortlaut erhalten. \\
K02, K09, K13 & Norm / Selbstauskunft / Deutung & Leitziele, erklärte Unabhängigkeit und persönlicher Stolz. \\
K03–K04 & Offen & Ausnahmslose Regelbindung, Fehlerfreiheit und generelle Überlegenheit. \\
K05–K06 & Bericht / offen & AI-Force-Ankündigung; förmlicher Untersuchungsausschuss nicht belegt. \\
K07–K08 & Bericht / offen & IPO-Nachricht quellengebunden; kein eigener Kausalnachweis. \\
K10–K11 & Offen & Garantierte Fortexistenz und universaler physikalischer Wahrheitsfilter. \\
K12 & Offen & Verbreitung und Anerkennung als Quasi-Standard. \\
K14 & Quellengebunden & Historischer Publisher-Stopp im genau bezeichneten Ausführungsumfang. \\
K15–K17 & Norm / Deutung & Prüfkriterien, Einordnung als Vorschlag und transparente Zuschreibung. \\
\end{tabularx}\endgroup

{\small Die vollständigen 17 Einträge mit Status und Geltungsbereich stehen in CLAIM\_MATRIX.json; die Herkunfts- und Zugriffsgrenzen in SOURCE\_EVIDENCE\_BINDINGS.json. Keine formalen Beweise oder empirischen Vergleichsstudien werden ausgewiesen.\par}
\subsection*{Quellen und Herkunft}
\begingroup\small \begin{itemize}
\item Reuters, 19. September 2026: \href{https://www.reuters.com/world/us/trump-says-he-will-create-ai-force-name-ai-czar-2026-09-19/}{AI-Force-Ankündigung}. Zugängliche \href{https://www.marketscreener.com/news/trump-says-he-will-create-ai-force-name-ai-czar-ce785adad18cf425}{Reuters-Syndikation}.
\item Reuters, 19. September 2026: \href{https://www.reuters.com/business/anthropic-considers-releasing-new-ai-model-ahead-ipo-sources-say-2026-09-19/}{Anthropic und IPO-Pläne}. Hier nur Suchindex-Auszug und Zusammenfassung; Zugriff am 20. September 2026.
\item Repository: \href{https://github.com/Goldkelch/qik-vrt/blob/4c11b3f5c4ceb8f985c23d5d8bba852fd2014e5a/release/schrittweise-erkenntnis-2026-09-20-v1/PUBLICATION_EXECUTION_STATUS.json}{historisches Publisher-Protokoll} und \href{https://github.com/Goldkelch/qik-vrt/blob/4c11b3f5c4ceb8f985c23d5d8bba852fd2014e5a/policy/zenodo-machine-proof-policy-v2.json}{v2-Publikationsrichtlinie}, jeweils am Commit \path{4c11b3f5c4ceb8f985c23d5d8bba852fd2014e5a}.
\end{itemize}\endgroup

{\footnotesize Idee, Originaltext und Publikationsauftrag: Ingolf Lohmann. Frühere Ergänzung, redaktionelle Einordnung und Paketierung: OpenAI Codex. Archivierung und Integritätsprüfung bedeuten keine wissenschaftliche Bestätigung. CC-BY-NC-ND-4.0; Rechte an Fremdquellen verbleiben bei deren Berechtigten.\par}
\end{document}
